\originalsection{向量空间、基与维数}
这一补充对应8.05原讲义Chapter 3全章的七节主题, 原作者Barton Zwiebach（2013年10月21日）并信用Sheldon Axler的线性代数讨论. 采用独立中文解释, 没有复制或打包Axler商业教材. 先从两能级系数$(c_1,c_2)$理解向量: 可加、可乘复数, 运算封闭, 有零矢和加法逆元. 这些代数运算满足交换、结合和分配律. 用复数作标量与用实数作标量会改变维数: $\mathbb C$自身在复数域维数1, 在实数域基为$1,\ii$, 维数2. 因而说“二维态”必须说明采用复振幅的维数.

一组矢\textbf{张成}空间, 指每个矢可写为它们的有限线性组合; \textbf{线性独立}指只有全零系数组合得零. 同时满足两者的是基, 系数展开唯一. 若有$n$个基矢, 维数为$n$; 矩阵空间的标准基是每次只在一格放1的$E_{ij}$, 维数为行数乘列数. 多项式空间包含任意长独立列$1,x,\ldots,x^N$, 因而无限维. 量子平方可积函数还需要范数极限来容纳无限展开, 这就是Hilbert空间比纯代数向量空间多的完整性条件.

子空间必须含零且对加法和倍乘封闭. $U+W=\{u+w\}$; 若交集只有零, 每个和分解唯一, 记$U\oplus W$. 正交直和另要求两子空间内任意两矢内积0. 代数直和不自动正交, 例如平面中的斜轴给唯一坐标却不是正交坐标.
\originalsection{线性映射、矩阵与本征矢}
$T:V\to W$满足$T(au+bv)=aTu+bTv$时称线性映射, $\ker T$是被映为0的矢, $\operatorname{im}T$是全部输出. \textbf{单射}指$Tu=Tv$必有$u=v$, 即不同输入不合并为同一输出; \textbf{满射}指每个$w\in W$都有$v\in V$使$Tv=w$, 即输出覆盖整个目标空间. 对线性映射, 检验单射只需解$Tv=0$并确认只有零解, 检验满射则确认$\operatorname{im}T=W$. 有限维中取零核基$k_1,\ldots,k_r$, 扩成全空间基$k_1,\ldots,k_r,v_1,\ldots,v_s$. 后面$s$个像张成像空间; 若$\sum a_jTv_j=0$, 则$\sum a_jv_j$落在零核, 与扩充基独立性矛盾, 故它们独立. 所以$\dim V=\dim\ker T+\dim\operatorname{im}T$. 相同有限维的输入输出间, 单射、满射、可逆等价; 无限维左移序列是满射而非单射, 不能照搬.

对$T:V\to V$, 选基后$Tv_j=\sum_iT_{ij}v_i$, 第$j$列表示第$j$基矢的输出. 复合对应矩阵乘法, 恒等对应单位阵. 对$n\times n$方阵, \textbf{迹}是对角元素之和, \textbf{行列式}是带排列符号的乘积和:
\[
 \Tr T=\sum_{i=1}^nT_{ii},\qquad
 \det T=\sum_{\sigma\in S_n}(-1)^{N_\sigma}\prod_{i=1}^nT_{i,\sigma(i)}.
\]
这里$S_n$为$1,\ldots,n$的全部排列, $N_\sigma$是满足$i<j$而$\sigma(i)>\sigma(j)$的逆序对数; 例如二维行列式为$T_{11}T_{22}-T_{12}T_{21}$. 实际计算可用一行加另一行的倍数消元, 这不改行列式; 每次交换两行变号, 上三角阵的行列式为对角元素之积. \textbf{特征多项式}定义为$p_T(z)=\det(zI-T)$, 可按行列式计算并整理成$z$的多项式. 换基列矩阵$C$满足$[v]_{\rm old}=C[v]_{\rm new}$, 算符新矩阵$C^{-1}TC$, 其迹、行列式和特征多项式不变. 这说明一个矩阵表格依基, 其物理作用不依基.

$Tu=\lambda u,u\ne0$定义本征对. 求本征值先解$p_T(\lambda)=0$, 再解$(T-\lambda I)u=0$找非零本征矢; 有对角本征基时, 迹与行列式分别给本征值之和与积, 可用来核对计算. 不同本征值的本征矢独立: 若最短非平凡依赖$\sum_{j=1}^na_ju_j=0$, 作用$T-\lambda_nI$消去最后一项, 剩下更短非平凡依赖, 矛盾. 复空间的特征多项式有根保证存在一个本征矢; 实矩阵如二维旋转可无实本征值. 有$n$个独立本征矢才可对角化, 存在一个本征矢远远不够. 例如$\begin{pmatrix}1&1\\0&1\end{pmatrix}$的$(T-I)(x,y)=0$只给$y=0$, 本征矢仅为$(x,0)$且$x\ne0$, 故本征空间一维, 不能对角化.
\originalsection{内积、正交化与投影}
内积$(u,v)$对第二变量线性、对第一变量共轭线性, 共轭交换, $\norm v^2=(v,v)>0$对非零$v$成立. 复空间若不共轭,$(1,\ii)$的“平方和”会为0而非正长度. Schwarz不等式来自将$u$减去沿$v$的最佳倍数; 三角不等式随之成立, 作业3有完整推导.

给独立列$v_1,\ldots,v_n$, Gram--Schmidt先$e_1=v_1/\norm{v_1}$, 再$w_j=v_j-\sum_{i<j}e_i(e_i,v_j)$, $e_j=w_j/\norm{w_j}$. 正交性由逐项内积核验; 独立性保证$w_j\ne0$. 它不会改变每一步已张成的子空间. 对子空间$U$的正交基,$P_Uv=\sum_je_j(e_j,v)$, 自然有$P_U^2=P_U=P_U^\dagger$. $v-P_Uv$与每个$e_j$正交, 对有限维$U$因而$V=U\oplus U^\perp$; 一般Hilbert空间中需$U$闭, 投影展开按范数收敛理解, 分解唯一. 以两部分的Pythagoras式证明$P_Uv$是唯一最佳近似, 详见作业4的函数逼近应用. 单有$P^2=P$只是斜投影, 还要Hermitian才能作为正交测量投影.
\originalsection{对偶、伴随与幺正}
线性泛函$\varphi$把矢映成复数, 且保持线性. 正交基中设$\varphi(e_j)=a_j$, 则$\varphi(v)=\sum a_jv_j=(u,v)$, 唯一$u=\sum a_j^*e_j$. 这就是有限维Riesz对应, 在Hilbert空间需要泛函连续. bra是这种泛函, 其系数恰为ket系数的复共轭.

伴随由$(u,Tv)=(T^\dagger u,v)$定义. 正交基中矩阵为共轭转置, 因为逐基内积给$(T^\dagger)_{ji}=T_{ij}^*$. 可核$(ST)^\dagger=T^\dagger S^\dagger$, 和的伴随逐项相加, 常数取共轭. Hermitian满足$T=T^\dagger$, 本征值实与异值本征矢正交的证明在下文给出谱分解. 有限维幺正满足$U^\dagger U=UU^\dagger=I$并保持所有内积; 反过来对每对矢保持内积便有$U^\dagger U=I$, 同有限维时单射亦满射, 再得$UU^\dagger=I$; 无限维只保内积的非满射等距嵌入不称幺正. 正交归一基经幺正作用仍正交归一. 幺正本征值模1: 对$Uu=zu$比较范数得$|z|=1$.

任意有限维复矩阵还有Schur上三角化: 选一个单位本征矢作第一基矢, 第一列除顶格外为0, 下方余块再递归酉换基, 得上三角. Hermitian上三角必须同时下三角, 所以对角. 更一般normal矩阵满足$TT^\dagger=T^\dagger T$, 上三角normal的第一行除对角元素外须为0（比较第一格对角的行范数与列范数）, 归纳仍对角. 这解释何种算符拥有完备正交本征基, 不是把所有矩阵都假设为可测观测量.

\originalsection{把代数规则补到可复算的算符语言}
Hermitian矩阵对实倍乘封闭, 对复倍乘不封闭（$\ii I$不Hermitian）; 它们构成实向量空间, $n\times n$时维数$n^2$, 与全部复矩阵的复维数$n^2$要区分. 无限序列空间含任意长的标准独立列, 函数空间亦有多项式独立列; 这说明无限维而非声称有限和能表示每个函数. 任一独立列不长于有限生成列: 逐次把独立矢代入生成式, 用一个非零系数替换对应生成矢, 保持生成; 若独立列更长, 最后一矢便落在前面独立矢张成空间, 矛盾. 由此所有有限基长度相同.

线性映射的和、倍乘与复合仍线性; 复合通常不交换. 在多项式上$D=\partial_x,M=x$满足$(DM-MD)f=f$. 序列的右移$R(a_0,a_1,\ldots)=(0,a_0,a_1,\ldots)$单射非满射, 左移$L(a_0,a_1,\ldots)=(a_1,a_2,\ldots)$满射非单射, $LR=I$但$RL\ne I$. 这提供了无限维不能凭单侧逆推出双侧逆的具体检查.

子空间$W$不变指$TW\subset W$, 才能将$T$限制为$W$上的算符. 本征空间$\ker(T-\lambda I)$是不变子空间, 包含零矢; 本征矢则必须非零. Jordan例中本征值1的代数重数2、几何重数1, 因而这两个“重数”不能互换.

正交投影在$U$上为恒等、在$U^\perp$上为零, 本征值只有0、1; 取适应该分解的有限正交基, 得$\operatorname{rank}P=\Tr P=\dim U$. $U=V$时$P=I$可逆, 只有真子空间投影才不可逆. 理想测量的子空间概率是$\norm{P\psi}^2$, 条件态为$P\psi/\norm{P\psi}$, 若概率0则该条件态未定义.

对固定$u$, 泛函$v\mapsto(u,Tv)$由Riesz对应唯一矢$T^\dagger u$, 故伴随存在唯一; 泛函对$u$共轭线性而Riesz对应也共轭, 两者抵消使$T^\dagger$对$u$线性. 定义立即给$(T^\dagger)^\dagger=T$, $(cT)^\dagger=c^*T^\dagger$及逆序乘积. 非正交基的Gram矩阵$G_{ij}=(v_i,v_j)$使内积为$x^\dagger Gy$, 因而伴随矩阵是$G^{-1}T^*G$, 其中$T^*$为共轭转置, 不能无条件用$T^*$代替伴随.

复空间若$(v,Av)=0$对全体$v$, 令$B(u,v)=(u,Av)$. 先用$u+v$扣除两个对角项得$B(u,v)+B(v,u)=0$, 再用$u+\ii v$得$\ii B(u,v)-\ii B(v,u)=0$, 故所有$B(u,v)=0$, $A=0$. 若全态$(v,Tv)$实, 则$(v,(T-T^\dagger)v)=0$, 上式推出$T=T^\dagger$; 反向由共轭交换立即成立. 实空间探针不能得到第二个式子, 这正是实反对称矩阵零二次型例的原因.

满射线性算符保持范数, 由复极化恢复全部内积, 因而幺正; 有限维中范数保持已保证单射并进而满射. 若$v_j=Uu_j$为两正交基, 则$U$在两基中的矩阵相同: $(v_i,Uv_j)=(Uu_i,U^2u_j)=(u_i,Uu_j)$. 这是同一个算符在特定两基的性质, 不是任意换基都不改变矩阵.


范数的三条性质为正定、$\norm{cv}=|c|\norm v$及三角不等式. 完备指每个Cauchy序列（任意$\epsilon>0$后项两两距离小于$\epsilon$）有空间内的范数极限, 有限维自动完备而一般函数子空间未必. 核与像对线性组合封闭且含零, 故确为子空间. 两个双侧逆若存在,$S=S(TQ)=(ST)Q=Q$, 所以逆唯一. 有限生成列删除冗余项可成基, 独立列持续添加外部矢直至张成亦成基. 有限维$T-\lambda I$不单射等价行列式0, 将本征空间和特征多项式联系起来.

Riesz表示唯一性也直接可核: 若$(u,v)=(u',v)$对全体$v$, 取$v=u-u'$得$\norm{u-u'}^2=0$. 伴随逆序证明为$(u,STv)=(S^\dagger u,Tv)=(T^\dagger S^\dagger u,v)$, 唯一性给结论. 例如正交基中$T=\begin{pmatrix}0&2&\ii\\1&-\ii&0\\3\ii&1&7\end{pmatrix}$, 则$T^\dagger=\begin{pmatrix}0&1&-3\ii\\2&\ii&1\\-\ii&0&7\end{pmatrix}$; 每格为对称位置的复共轭. 对无界算符仍必须保留后文所说的定义域条件, 不能以有限矩阵论证替代.

